\section{Conclusion}
\label{sec:conclusion}

We present a systematic calibration audit of ECG foundation models: six encoders---five ECG
foundation models (MERL, ST-MEM, HuBERT-ECG, ECGFM-KED, ECG-FM) and a supervised S4D
baseline---across four datasets, with 23 of 24 model--dataset cells evaluated
(ECGFM-KED~$\times$~CODE-15\% excluded for a partial checkpoint), over 18{,}877 ECG records
(2{,}000 PTB-XL; 6{,}877 CPSC2018; 8{,}000 CODE-15\%; 2{,}000 MIMIC-IV-ECG) under
patient-level leakage-controlled splits with bootstrap confidence intervals.

\paragraph{Main findings.}
Four findings stand out. First, a consistent \textbf{discrimination--calibration tradeoff}
on PTB-XL: the five foundation models span AUROC 0.769--0.838 but all trail the supervised
S4D baseline's calibration (ECE 0.064), exceeding it by 0.015--0.072. ST-MEM achieves the
best discrimination (AUROC 0.838) but poor calibration (ECE 0.124), and HuBERT-ECG is the
worst calibrated (ECE 0.136)---a gap invisible to AUROC reporting but clinically significant.
Second, this tradeoff \textbf{disappears on arrhythmia datasets}: on CPSC2018
(ECE 0.023--0.051), CODE-15\% (ECE 0.010--0.032), and MIMIC-IV-ECG (ECE 0.043--0.106),
encoders are comparatively well-calibrated and ST-MEM leads discrimination. Task
structure---correlated multi-label vs.\ lower-co-occurrence endpoints---is the mediating factor. Third, \textbf{post-hoc
isotonic recalibration closes 43--71\% of the ECE gap} on PTB-XL (ST-MEM: $0.124\to0.044$;
HuBERT-ECG: $0.136\to0.040$) at only 0.5--1.7\% AUROC cost, and it reorders calibration-aware
model selection (ST-MEM moves from 5th to 1st on the $C(0.7)$ composite). Fourth,
\textbf{within-domain pretraining does not guarantee superiority}: ECG-FM, pretrained on
MIMIC-IV-ECG, achieves only AUROC 0.771 on MIMIC-IV-ECG---behind MERL (0.861) and ST-MEM
(0.854) in all 10 resplits, and no better than the supervised S4D baseline (0.793).
Cross-dataset ECE varies by up to ${\sim}6\times$ for
the same model, mandating per-deployment calibration audits.

\paragraph{A practical recipe.}
Beyond the audit we contribute \textbf{RankSafe-Cal}, which recovers most of isotonic's
calibration gain (PTB-XL macro ECE $0.097\!\to\!0.054$; on MIMIC it beats unconstrained
isotonic) at roughly $6\times$ lower AUROC cost by certifying out-of-sample ranking loss
within a tolerance $\epsilon$. Cross-dataset transfer experiments show recalibration ports
successfully in only 9--42\% of endpoint/model pairs, so local re-calibration remains
necessary. Split-conformal prediction attains its $90\%$ coverage guarantee as a complement
to calibration, and decision-curve analysis confirms positive net benefit for MI and AF
endpoints. All headline conclusions hold across 10 patient-level resplits.

\paragraph{Clinical implication.}
Practitioners deploying ECG foundation models for multi-label risk stratification should
mandate calibration auditing alongside AUROC benchmarking. Raw foundation-model
probabilities can be unusable for threshold-based clinical decisions even when AUROC is
excellent; post-hoc recalibration with as few as 63 examples eliminates most of the
calibration gap. The large cross-dataset ECE variation (e.g., ST-MEM ECE 0.028 on
CODE-15\% vs.\ 0.124 on PTB-XL) means calibration audits must be conducted on locally
representative patient cohorts---not extrapolated from published benchmark datasets.

\paragraph{Pretraining objective signal.}
A graded signature emerges: MERL (contrastive InfoNCE) shows the strongest underconfidence
($\text{SCE}=-0.003$), masked-prediction models (ST-MEM, HuBERT-ECG) are near-neutral, and
knowledge-enhanced, hybrid, and supervised models are mildly overconfident
($+0.007$--$+0.008$), consistent with the documented underconfidence of
contrastive FMs. Feature geometry analysis shows that MERL's contrastive objective
produces extreme feature spread (std of L2 norms: 34.8) relative to masked-prediction
and supervised models. Notably, MERL achieves the best performance on MIMIC-IV-ECG
(AUROC 0.866, ECE 0.033) despite not being pretrained on that cohort.

\paragraph{Open problems.}
Three directions remain for future work. First, \textbf{fine-tuning calibration}: this
study evaluates frozen linear probes; full fine-tuning may alter calibration in complex
ways that require re-auditing. Second, \textbf{larger MIMIC-IV-ECG evaluation}: this
benchmark used 2{,}000 records; the full 800{,}035-record corpus with structured ICD-based
labels (rather than parsed machine-measurement reports) would further strengthen the MIMIC
conclusions. Third, \textbf{objective-controlled pretraining}: the graded calibration
signature across contrastive, masked-prediction, knowledge-enhanced, and generative
objectives motivates controlled experiments that vary only the pretraining objective to
establish causality.

\paragraph{Call to action.}
The \texttt{ecg-audit} package and all benchmark scripts are ready to run against new
checkpoints without modification. We invite the ECG AI community to submit calibration
results to the leaderboard, enabling cumulative evidence on which models can be safely
deployed as probability estimators in clinical cardiac decision support.
